\documentclass[11pt,a4paper]{article}

\usepackage[margin=0.85in]{geometry}
\usepackage{graphicx}
\usepackage{booktabs}
\usepackage{array}
\usepackage{xcolor}
\usepackage{colortbl}
\usepackage{hyperref}
\usepackage{enumitem}
\usepackage{longtable}

\hypersetup{colorlinks=true, linkcolor=black, urlcolor=blue}

\definecolor{banditblue}{HTML}{dce7f5}
\definecolor{rowalt}{HTML}{f3f3f3}

\setlength{\parskip}{0.45em}
\setlength{\parindent}{0pt}

\title{\vspace{-1.5em}\textbf{Project Management Report} \\ \large Distributed Financial Transaction Processor\vspace{-0.3em}}
\author{Sampath Pranay Beela, Rahul Chinya Jagadeesha \\ \small CS6650 --- Building Scalable Distributed Systems}
\date{\vspace{-0.5em}\small Spring 2026}

\begin{document}
\maketitle
\vspace{-1.8em}

\section*{1. Initial design vs.\ final state}

We began with a one-page sketch of the system: an HTTP API in front, a worker pool in the back, and a queue in between. The sketch answered ``what shape is this going to be'' but nothing else. The final system is a fully cloud-deployed, Terraform-provisioned, Locust-benchmarked, four-experiment-validated pipeline running on real AWS infrastructure. Table~\ref{tab:delta} captures the journey from intent to outcome.

\begin{table}[h]
\centering
\small
\renewcommand{\arraystretch}{1.22}
\rowcolors{2}{white}{rowalt}
\begin{tabular}{>{\raggedright\arraybackslash}p{4.3cm} >{\raggedright\arraybackslash}p{4.9cm} >{\raggedright\arraybackslash}p{4.9cm}}
\toprule
\rowcolor{banditblue}
\textbf{Aspect} & \textbf{Initial intent (Week 1)} & \textbf{Final state (Week 8)} \\
\midrule
Architecture & API + worker + queue + DB & ALB $\rightarrow$ ECS Fargate api-svc $\rightarrow$ SQS (+ DLQ) $\rightarrow$ ECS Fargate worker-svc $\rightarrow$ DynamoDB (3 tables) + CloudWatch \\
Deployment target & ``we'll figure out AWS later'' & Terraform-managed stack in \texttt{us-west-2}, preflight-gated under LabRole \\
Concurrency control & optimistic only & optimistic + pessimistic, swap by env var, compared head-to-head \\
Correctness story & ``we'll make it atomic'' & single \texttt{TransactWriteItems} for debit + credit + status; crash-test-verified \\
Observability & console logs & CloudWatch log groups for both services + six Prometheus-style worker counters + JSON \texttt{METRICS} log lines every 30~s \\
Load testing & ``some Locust script'' & two user classes, four experiment runners, auto-rendered chart pipeline \\
Reproducibility & manual steps & every task a preflight-gated \texttt{make} target; every experiment a shell script that reseeds, runs, drains, verifies \\
\bottomrule
\end{tabular}
\caption{Where the project started vs.\ where it ended.}
\label{tab:delta}
\end{table}

Three things changed the most between ``initial intent'' and ``final state.'' First, we committed early to a \textbf{cloud-only} workflow rather than a local-compose-then-cloud workflow --- this made the middle of the project harder but produced experimental numbers we could actually defend. Second, we decided to wrap the debit, credit, and status update in a single \texttt{TransactWriteItems} rather than three separate writes, which turned out to be the single design decision that the crash-recovery experiment proved correct. Third, we added a second concurrency-control strategy (pessimistic locking) specifically so Experiment~3 could answer a real question, not just measure our one implementation.

\section*{2. How we broke the problem down}

We used the eight-week timeline in the project implementation plan as a living roadmap. Each of the four two-week phases owned a distinct set of deliverables with explicit dependencies so work could flow between phases without stalling. Table~\ref{tab:phases} is the phase breakdown we actually followed.

\begin{table}[h]
\centering
\small
\renewcommand{\arraystretch}{1.25}
\rowcolors{2}{white}{rowalt}
\begin{tabular}{>{\raggedright\arraybackslash}p{2.2cm} >{\raggedright\arraybackslash}p{7.1cm} >{\raggedright\arraybackslash}p{5.0cm}}
\toprule
\rowcolor{banditblue}
\textbf{Phase} & \textbf{Deliverables} & \textbf{Success criteria} \\
\midrule
Weeks 1--2 \newline \textit{Foundations} & Architecture finalization; repo + service scaffolding; DynamoDB schema; core transfer workflow; API-side idempotency; SQS producer and consumer & One transfer can be submitted via HTTP and is applied exactly once to DynamoDB \\
Weeks 3--4 \newline \textit{Cloud + first load} & Dockerfiles + ECR; Terraform (ECR, ECS, ALB, SQS, Dynamo, CloudWatch); cloud smoke test; Locust harness (two workload classes); \textbf{Exp~1} (hot-account storm); \textbf{Exp~2} (worker crash recovery) & System runs on AWS end-to-end; two experiments produce reproducible numbers \\
Weeks 5--6 \newline \textit{Comparison + scale} & Pessimistic locking implementation; mode-swap harness; \textbf{Exp~3} (opt vs.\ pess on cloud); \textbf{Exp~4} (horizontal scaling 1:1, 2:2, 4:4) & Both modes characterized; scaling curve produced; correctness held at every config \\
Weeks 7--8 \newline \textit{Analysis + delivery} & Chart generator; summary-markdown generator; final report (LaTeX, 24 pages with embedded AWS screenshots); experiments report (5 pages); lessons-learned reflections; demo prep & Publication-ready artifacts produced; demo runnable end-to-end from a clean shell \\
\bottomrule
\end{tabular}
\caption{Four phases, each with a clear ``we move on when this is true'' gate.}
\label{tab:phases}
\end{table}

Within each phase, we worked in short commit-sized units: each commit landed one self-contained piece (a handler, a script, a Terraform resource, an experiment runner) behind a feature flag or behind a preflight check, so neither of us ever blocked the other's work. This mattered the most during Weeks~3--4, when one of us was writing Terraform while the other was writing the crash-recovery experiment against the output of that Terraform.

\section*{3. Who worked on what}

The team is two people. The work split came directly from the project implementation plan and held without revision through all eight weeks. Table~\ref{tab:ownership} is the per-item RACI view.

\begin{table}[h]
\centering
\small
\renewcommand{\arraystretch}{1.22}
\rowcolors{2}{white}{rowalt}
\begin{tabular}{>{\raggedright\arraybackslash}p{5.6cm} >{\raggedright\arraybackslash}p{4.0cm} >{\raggedright\arraybackslash}p{4.9cm}}
\toprule
\rowcolor{banditblue}
\textbf{Work item} & \textbf{Owner} & \textbf{Artifact} \\
\midrule
Project scope and architecture finalization & Pranay \& Rahul & \texttt{arch.md}, design discussions \\
Repository and service scaffolding          & Pranay          & \texttt{api-svc/}, \texttt{worker-svc/} \\
DynamoDB schema design                      & Rahul           & \texttt{accounts}, \texttt{transactions}, \texttt{account\_locks} tables \\
Core transfer workflow                      & Pranay \& Rahul & API handler (Pranay), worker commit primitives (Rahul) \\
API-layer idempotency                       & Rahul           & \texttt{PutTransactionIfNotExists} + handler short-circuit \\
SQS integration                             & Pranay          & \texttt{api-svc/queue}, \texttt{worker-svc/queue} \\
Local correctness verification              & Rahul           & \texttt{load-tests/verify\_balances.py} \\
Containerization                            & Pranay          & multi-stage \texttt{Dockerfile}s, \texttt{scripts/push\_images.sh} \\
Terraform infrastructure                    & Rahul           & \texttt{infra/} \\
Cloud smoke testing                         & Pranay \& Rahul & \texttt{scripts/smoke\_test.sh}, \texttt{make cloud-smoke} \\
Load testing framework                      & Pranay          & \texttt{load-tests/locustfile.py}, two user classes \\
Experiment 1 --- hot-account storm          & Pranay          & \texttt{make cloud-test-load-hot}, results CSVs \\
Experiment 2 --- worker crash recovery      & Rahul           & \texttt{scripts/crash\_recovery\_test.sh} \\
Pessimistic locking implementation          & Rahul           & \texttt{worker-svc/processor} pessimistic path, \texttt{account\_locks} primitives \\
Experiment 3 --- opt vs.\ pess              & Pranay \& Rahul & \texttt{scripts/compare\_locking\_modes.sh} \\
Experiment 4 --- horizontal scaling         & Pranay \& Rahul & \texttt{scripts/scale\_experiment.sh} \\
Data analysis and chart generation          & Pranay          & \texttt{scripts/render\_charts.py}, \texttt{summarize\_locust\_results.py} \\
Final report writing                        & Rahul           & \texttt{load-tests/final-report.tex} \\
Final demo preparation                      & Pranay \& Rahul & demo script + recorded walkthrough \\
\bottomrule
\end{tabular}
\caption{Per-item ownership. ``Pranay \& Rahul'' items were paired work where both of us contributed meaningfully.}
\label{tab:ownership}
\end{table}

We kept the split clean by designing around interfaces: the API handler (Pranay) called into DynamoDB accessors (Rahul) via a typed Go interface, and the two concurrency-control paths (Rahul's pessimistic, Pranay-integrated with the existing optimistic) lived behind the same \texttt{Processor.attemptTransfer} method. Neither of us had to wait for the other to finish a module before we could start ours.

\section*{4. Problems encountered and how they were resolved}

We kept an informal issue log as we went. The entries that mattered most to the final state are summarized in Table~\ref{tab:issues}. Each entry is a real obstacle that cost at least one session to diagnose.

\begin{table}[h]
\centering
\small
\renewcommand{\arraystretch}{1.2}
\rowcolors{2}{white}{rowalt}
\begin{tabular}{>{\raggedright\arraybackslash}p{3.7cm} >{\raggedright\arraybackslash}p{5.0cm} >{\raggedright\arraybackslash}p{5.7cm}}
\toprule
\rowcolor{banditblue}
\textbf{Problem} & \textbf{Root cause} & \textbf{Resolution} \\
\midrule
ECS \texttt{CannotPullContainerError} on first deploy & Images built on Apple Silicon (\texttt{arm64}); Fargate defaults to \texttt{linux/amd64} & Switched \texttt{scripts/push\_images.sh} to \texttt{docker buildx build --platform linux/amd64} \\
\texttt{terraform apply} fails with ``no matching EC2 VPC found'' & No default VPC in \texttt{us-west-2} in this AWS account & \texttt{aws ec2 create-default-vpc --region us-west-2} once; unblocks all applies \\
\texttt{LabRole} not present in the lab account & Our Northeastern SSO account is not AWS Academy-shaped & Wrote \texttt{scripts/bootstrap\_labrole.sh} to create the role and attach the four managed policies; tightened the preflight check to assert the role is ECS-assumable rather than assert the caller \textit{is} LabRole \\
\texttt{aws iam create-role} failing with shell-quoting errors & \texttt{zsh} brace expansion eating multiline JSON arguments & Moved trust policy to \texttt{infra/labrole-trust.json}; every \texttt{aws iam} call now uses \texttt{file://} \\
\texttt{date +\%s\%3N} producing literal \texttt{...3N} on macOS & macOS \texttt{date} has no \texttt{\%N} format & Replaced with \texttt{python3 -c 'import time; print(int(time.time()*1000))'} in both sweep scripts \\
Between-iteration reseed hitting \texttt{TransactionConflictException} & In-flight worker \texttt{TransactWriteItems} hold a short lock on the item being \texttt{PutItem}'d by reseed & Added SQS purge + 65-second visibility drain before reseed; added 8-attempt exponential backoff retry in the seed script for \texttt{TransactionConflictException} \\
Transient balance mismatch during scaling sweep & Verifier ran before SQS queue fully drained; verifier used eventually consistent \texttt{Scan} & Added drain loop (visible + in-flight = 0) with 15-second settle buffer; switched \texttt{verify\_balances.py} to \texttt{ConsistentRead=True} \\
One verify failure aborting whole sweep & \texttt{set -e} treated the transient mismatch as fatal & Sweep scripts now log a warning, write the verify text to a file, and continue to the next iteration \\
Locust venv import error (\texttt{zope.event}) & Stale dependency in shared venv & Rebuilt \texttt{/tmp/locust-venv} from scratch; added \texttt{boto3} so \texttt{verify\_balances.py} runs in the same interpreter \\
\bottomrule
\end{tabular}
\caption{Nine issues we hit, each with root cause and resolution. Every one of these is discussed in more depth in \texttt{final-report.pdf}.}
\label{tab:issues}
\end{table}

The issues cluster into three categories. The \textbf{tooling} category (Docker cross-build, shell quoting, macOS \texttt{date}, Locust venv) were one-time environment mismatches --- annoying but solvable with one-liner changes. The \textbf{cloud-identity} category (LabRole, VPC) was about our account not matching the assumptions baked into the assignment; fixing these taught us more about AWS IAM trust policies than the rest of the project combined. The \textbf{correctness-under-concurrency} category (reseed conflicts, drain waits, \texttt{ConsistentRead}) was the most educational --- those were all places where the system was doing the right thing but our \textit{observation} of the system was wrong, and fixing them required us to separate ``is the system correct'' from ``can I see that it is correct.''

\section*{5. Timeline of commits and milestones}

We did not keep formal sprint retros, but every pull-request-sized unit of work maps to a phase milestone. At a high level the commit cadence looked like this: Weeks~1--2 was heavy on scaffolding commits and the first end-to-end transfer; Weeks~3--4 was dominated by Terraform + cloud-first deployment and the first two experiments; Weeks~5--6 was pessimistic locking plus the two comparison experiments; Weeks~7--8 was documentation, screenshots, chart rendering, and reports. The repository's commit history reflects that cadence --- not the number of commits (the exact count is not the right metric), but the shape of them: feature commits early, infrastructure commits in the middle, refinement and analysis commits at the end.

\section*{6. What we would do differently next time}

Three process-level changes we would make if we started over:

\begin{enumerate}[leftmargin=*, itemsep=0.15em]
    \item \textbf{Build the verification harness with the same care as the service.} Our transient-mismatch scare during Experiment~4 came from the verifier using an eventually consistent Scan, not from the system losing money. If we had treated \texttt{verify\_balances.py} as first-class code from day one (with \texttt{ConsistentRead=True}, a drain gate, and proper error discrimination between ``mismatch'' and ``still in flight''), we would have saved ourselves an hour of debugging a non-bug.

    \item \textbf{Bootstrap the cloud identity story on day one.} We initially assumed a LabRole existed in our AWS account. It didn't. Creating LabRole ourselves (\texttt{scripts/bootstrap\_labrole.sh}) should have happened during Weeks~1--2, not discovered halfway through Weeks~3--4. The general lesson: verify every external assumption your stack makes \textit{before} you start depending on it.

    \item \textbf{Cross-platform build defaults early.} A single-line \texttt{docker buildx --platform linux/amd64} in the original \texttt{push\_images.sh} would have saved an entire deploy cycle during the first ECS rollout. We now treat architecture pinning as a day-one concern on any project that targets a cloud runtime.
\end{enumerate}

\section*{7. Final state}

Every deliverable in the project implementation plan is done. The system is deployable end-to-end via \texttt{make tf-apply \&\& make push-images \&\& make tf-apply \&\& make cloud-seed \&\& make cloud-smoke}, and tears down cleanly via \texttt{make tf-destroy}. Each experiment has its own \texttt{make} target that produces timestamped artifacts under \texttt{load-tests/experiments/}, including charts and a markdown summary. Four reports live under \texttt{load-tests/}: \texttt{final-report.pdf} (comprehensive individual report), \texttt{experiments-report.pdf} (5-page experiments summary), \texttt{lessons-learned.pdf} (individual reflection), and this project-management report. The full repository is self-contained and reproducible on any AWS account that has a LabRole-shaped role, with no elevated permissions required at apply time.

\end{document}
